\section{Semi-continuity}

In this section, we fix an integral noetherian base scheme $S$, a flat projective morphism $X \to S$ with integral fibers, and $f: X \dashrightarrow X$ a family of dominant rational maps over $S$.
For each $s \in S$, we denote by $X_s$ the fiber of $X$ over $s$, and by $f_s: X_s \dashrightarrow X_s$ the induced rational map on the fiber.

% \subsection{Semi-continuous functions on noetherian schemes}

%   Let $S$ be an integral noetherian scheme, and $\theta: S \to \bbR$ be a function.
%   We say that $f$ is \emph{lower semi-continuous} if for every $r \in \bbR$, the set $\{ s \in S : f(s) > r \}$ is open in $S$.

%   \begin{remark}
%     The limit of a sequence of lower semi-continuous functions may not be lower semi-continuous. 
%     \Yang{}
%   \end{remark}

%   \begin{lemma}
%     A function $\theta: S \to \bbR$ is lower semi-continuous if and only if the following condition holds: 
%     \begin{enumerate}
%       \item for every $\eta \in S$ and every specialization $\xi \in \overline{\{\eta\}}$, we have $\theta(\xi) \leq \theta(\eta)$;
%       \item for every $\eta \in S$ and $a < \theta(\eta)$, there is an open subset $U \subseteq \overline{\{\eta\}}$ such that $\theta(\xi) > a$ for every $\xi \in U$.
%     \end{enumerate}
%   \end{lemma}
%   \begin{proof}
%     \Yang{}
%   \end{proof}


%   \begin{lemma}
%     Let $\theta: S \to \bbR$ be a lower semi-continuous function.
%     Then 
%     \begin{enumerate}
%       \item $\theta$ is continuous with respect to the constructible topology on $S$;
%       \item if $\theta$ is discrete, then its image is finite.
%     \end{enumerate}
%   \end{lemma}
%   \begin{proof}
%     \Yang{}
%   \end{proof}


\subsection{Semi-continuity of mixed degrees}

  \begin{proposition}
    % Let $f: X \dashrightarrow X$ be a family of dominant rational maps over $S$.
    Fix a $\pi$-ample divisor $L$ on $X$ and integers $m_1, \ldots, m_d \in \bbZ_{\geq 0}$.
    Then the function 
    \[ s \mapsto (f_s^{m_1})^*L_s \cdot \cdots \cdot (f_s^{m_d})^*L_s \]
    is lower semi-continuous on $S$.
  \end{proposition}
  \begin{proof}
    Let $\Gamma_{f,\bfm} \subset X \times_S \cdots \times_S X$ be the closure of the image of the rational map $(f^{m_1}, \ldots, f^{m_d}): X \dashrightarrow X \times_S \cdots \times_S X$.
    By Gruson-Raynaud's flattening theorem (\Yang{}), there exists a blowup $S' \to S$ such that the strict transform $\Gamma'_{f,\bfm}$ of $\Gamma_{f,\bfm}$ is flat over $S'$.
    Let $X', L', f'$ be the base change of $X, L, f$ along $S' \to S$ respectively.
    For every $s' \in S'$ with image $s \in S$, the date $(X'_{s'}, L'_{s'}, f'_{s'})$ is the base change of $(X_s, L_s, f_s)$.
    Hence the semi-continuity on $S$ follows from the semi-continuity on $S'$.
    By replacing $S$ with $S'$, we may assume that $\Gamma_{f,\bfm}$ is flat over $S$.

    For every $s \in S$, let $\Gamma_{f_s,\bfm}$ be the closure of the image of the rational map $(f_s^{m_1}, \ldots, f_s^{m_d}): X_s \dashrightarrow X_s \times \cdots \times X_s$.
    Denote by $\pi_i: X \times_S \cdots \times_S X \to X$ the projection to the $i$-th factor.
    Then we have
    \[ (f_s^{m_1})^*L_s \cdot \cdots \cdot (f_s^{m_d})^*L_s = \pi_1^*L \cdot \cdots \cdot \pi_d^*L \cdot [\Gamma_{f_s,\bfm}]. \]

    Over the generic point $\eta$ of $S$, we have $(\Gamma_{f,\bfm})_\eta = \Gamma_{f_\eta,\bfm}$.
    Over a special point $s \in S$, $(\Gamma_{f,\bfm})_s$ may have several irreducible components, and $\Gamma_{f_s,\bfm}$ is one of the irreducible components of $(\Gamma_{f,\bfm})_s$.
    Since $\Gamma_{f,\bfm}$ is flat over $S$ and intersection numbers are constant in flat families, we have
    \begin{align*}
      (f_s^{m_1})^*L_s \cdot \cdots \cdot (f_s^{m_d})^*L_s & = \pi_1^*L \cdot \cdots \cdot \pi_d^*L \cdot [\Gamma_{f_s,\bfm}] \\
      & \leq \pi_1^*L \cdot \cdots \cdot \pi_d^*L \cdot [(\Gamma_{f,\bfm})_s] \\
      & = \pi_1^*L \cdot \cdots \cdot \pi_d^*L \cdot [(\Gamma_{f,\bfm})_\eta] \\
      & = \pi_1^*L \cdot \cdots \cdot \pi_d^*L \cdot [\Gamma_{f_\eta,\bfm}] \\
      & = (f_\eta^{m_1})^*L_\eta \cdot \cdots \cdot (f_\eta^{m_d})^*L_\eta.
    \end{align*}
    And since the generic fiber is (geometrically) irreducible, there exists a non-empty open subset $U \subseteq S$ such that for every $s \in U$, we have $(\Gamma_{f,\bfm})_s$ is irreducible and hence equal to $\Gamma_{f_s,\bfm}$.
    Then for every $s \in U$, the above inequalities are equalities.
    By \Yang{}, we are done.
    % \Yang{}
  \end{proof}

  \begin{example}
    Let $S = \Spec \bbZ$, $X = \bbP^2_{\bbZ}$, $L = \calO(1)$, and $f: X \dashrightarrow X$ is defined by
    \[ [x:y:z] \mapsto [xz:yz+2xy:z^2]. \]
    Then the function $s \mapsto \deg_1(f_s) = L_s \cdot (f_s^*L_s)$ is not constant on $S$.
    Indeed, we have $\deg_1(f_{(2)}) = 1$ and $\deg_1(f_{(p)}) = 2$ for every prime $p \neq 2$.
    
    Let $\Gamma \subset X \times_S X = \bbP^2_{\bbZ} \times \bbP^2_{\bbZ}$ be the graph of $f$.
    On $\bbP^2_{\bbZ} \times \bbP^2_{\bbZ}$, we use the coordinates $[x:y:z]$ and $[u:v:w]$ for the first and second factor respectively.
    Set $V = V(u(yz+2xy) - vxz, w(yz+2xy) - vz^2, wx - uz)$.
    Then $\Gamma$ is an irreducible component of $V$.
    Consider the first projection $\pi_1: V \to X$.
    Over the open subset $z \neq 0$ of $X$, $\pi_1$ is an isomorphism.
    When $x,y \neq 0$ and $z = 0$, $\pi_1$ is also an isomorphism.
    Over the point $[1:0:0]$, the fiber of $\pi_1$ is a line in the second factor defined by $w = 0$.
    Over the point $[0:1:0]$, the fiber of $\pi_1$ is the full fiber $\{[0:1:0]\} \times \bbP^2$.
    Hence $V = \Gamma \cup (\{[0:1:0]\} \times \bbP^2)$.
    We need to find a function vanishing on $\Gamma$ but not vanishing on $\{[0:1:0]\} \times \bbP^2$.
    Note that $(2u+w)y-vz$ satisfies this condition.
    Hence $\Gamma$ is given by 
    \[ \Gamma = V(u(yz+2xy) - vxz, w(yz+2xy) - vz^2, wx - uz, (2u+w)y-vz). \]

    Consider the first projection $\pi_1: \Gamma \to X$.
    Over any point $s \in \Spec \bbZ$ different from $(2)$, $\pi_1$ is an isomorphism over $\bbP^2_s \setminus \{[1:0:0],[0:1:0]\}$, and over this two points, the fiber of $\pi_1$ is a line.
    Hence $\pi_1$ is the blowup of $X_s$ at these two points.
    Over the point $(2)$, $\Gamma_{(2)}$ is given by 
    \begin{align*}
      \Gamma_{(2)} & =  V(uyz - vxz, wyz - vz^2, wx - uz, wy-vz) \\
        &= V((uy-vx)z, wx - uz, wy-vz) \\
        &= V(uy-vx, wx - uz, wy-vz) \cup V(z, w) \\
        &= \Delta \cup (\{z=0\} \times \{w=0\}).
    \end{align*}
    It has two irreducible components, one is the diagonal $\Delta$ and the other is $\{z=0\} \times \{w=0\}$.
    Among them, $\Delta$ is the graph of $f_{(2)}$ and $\{z=0\} \times \{w=0\}$ is the ``extra'' component.
    % \Yang{}
  \end{example}


\subsection[Semi-continuity of dynamical degree]{Semi-continuity of $\lambda_i(f)$}

  \begin{theorem}
    The function $s \mapsto \lambda_i(f_s)$ is lower semi-continuous on $S$.
  \end{theorem}
  \begin{proof}
    Fix a $\pi$-ample divisor $L$ on $X$.
    It determines a class $L_s \in \upN^1(\frakX_s)$ for every $s \in S$.

    Let $\eta$ be the generic point of $S$.
    For any $s \in S$, we have $\deg_i(f_s^m) \leq \deg_i(f_\eta^m)$ for every $m$ by the previous proposition, which implies $\lambda_i(f_s) \leq \lambda_i(f_\eta)$.

    Hence we only need to show that for every $\beta < \lambda_i(f_\eta)$, there is a nonempty open subset $U \subseteq S$ such that $\lambda_i(f_s) > \beta$ for every $s \in U$.
    \Yang{}


    \Yang{}
  \end{proof}

  \begin{example}
    Let $S = \Spec \bbZ$, $X = \bbP^2_{\bbZ}$, and $f: X \dashrightarrow X$ is defined by
    \[ [x:y:z] \mapsto [xy:xy-2z^2:yz+3z^2]. \]
    Then we have $\lambda_1(f_{(0)}) = 2$ and $\lambda_1(f_{(p)}) < 2$ for every prime $p > 2$.
    
    \Yang{}
  \end{example}
